\documentclass[a4paper,12pt]{article}
\usepackage[T2A]{fontenc}
\usepackage[utf8]{inputenc}
\usepackage[russian]{babel}
\usepackage{amsmath,amssymb,mathtools}
\usepackage{helvet}
\renewcommand{\familydefault}{\sfdefault}
\usepackage{icomma}
\usepackage[margin=18mm]{geometry}
\usepackage{microtype}
\usepackage{xcolor}
\usepackage{enumitem}
\usepackage{tabularx,booktabs}
\usepackage{parskip}
\usepackage{setspace}
\onehalfspacing
\definecolor{accent}{HTML}{287C7A}
\definecolor{darkaccent}{HTML}{184E4D}
\definecolor{textgray}{HTML}{253033}
\color{textgray}
\setlist[itemize]{leftmargin=1.5em,itemsep=.25em}
\setlist[enumerate]{leftmargin=1.8em,itemsep=.35em}
\begin{document}
\begin{titlepage}
\centering
\vspace*{30mm}
{\color{darkaccent}\Large\bfseries ЧЕК-ЛИСТ}\\[7mm]
{\Huge\bfseries Комплексное повторение\\[2mm]геометрии и статистики}\\[5mm]
{\Large Подготовка к МЦКО}\\[20mm]
{\large Екатерина Скелли}\\
\vfill
{\large 26.09.2026}\\[8mm]
{\small Лёвин Артём Александрович\\эксклюзивно для Екатерины Скелли}
\end{titlepage}

\section*{\color{darkaccent}1. Главная стратегия}
Если требуется найти отрезок или угол, сначала определите, в какой треугольник входит искомый элемент. Особенно полезны прямоугольные и равнобедренные треугольники.
\begin{enumerate}
\item Отметьте искомый элемент.
\item Найдите подходящий треугольник или угловую конфигурацию.
\item Узнайте тип: равнобедренный, прямоугольный, параллельные прямые, окружность.
\item Примените теорему.
\item При необходимости выразите одну величину двумя способами и составьте уравнение.
\item Проверьте ответ.
\end{enumerate}

\section*{\color{darkaccent}2. Неравенство треугольника}
Для сторон треугольника
\[
a+b>c,\qquad a+c>b,\qquad b+c>a.
\]
Для неизвестной третьей стороны удобно:
\[
\boxed{|a-b|<x<a+b}.
\]
Для сторон $7$ и $11$:
\[
4<x<18,
\]
поэтому положительные целые значения: $5,6,\ldots,17$; всего $13$.

\section*{\color{darkaccent}3. Равнобедренный треугольник}
Если $AB=BC$, то углы при основании равны. Медиана, проведённая из вершины к основанию, одновременно является высотой и биссектрисой.

\section*{\color{darkaccent}4. Прямоугольный треугольник}
Острые углы в сумме дают $90^\circ$.
Катет, лежащий против угла $30^\circ$, равен половине гипотенузы.
Если один острый угол равен $45^\circ$, второй также равен $45^\circ$, поэтому катеты равны.
Медиана, проведённая к гипотенузе, равна половине гипотенузы:
\[
BM=AM=CM=\frac{AC}{2}.
\]

\section*{\color{darkaccent}5. Опорная задача}
Пусть $AB=BC$, $\angle ABC=120^\circ$, $BM$ --- медиана к основанию $AC$. На луче $BM$ за $M$ выбрана точка $F$, $AF\perp AB$, $MF=72$. Обозначим $AB=x$.
Из свойств равнобедренного треугольника $\angle ABM=60^\circ$, поэтому в прямоугольном $ABM$
\[
BM=\frac{x}{2}.
\]
В прямоугольном $ABF$ угол $F$ равен $30^\circ$, поэтому
\[
BF=2x.
\]
Но
\[
BF=BM+MF=\frac{x}{2}+72.
\]
Значит,
\[
2x=\frac{x}{2}+72,\qquad 4x=x+144,\qquad x=48.
\]

\section*{\color{darkaccent}6. Углы}
Смежные углы:
\[
\alpha+\beta=180^\circ.
\]
Если углы относятся как $a:b$, удобно обозначить их $ax$ и $bx$.
Биссектриса делит угол пополам.
Биссектрисы двух смежных углов всегда перпендикулярны.

\section*{\color{darkaccent}7. Параллельные прямые}
При двух параллельных прямых и секущей:
\begin{itemize}
\item соответственные углы равны;
\item накрест лежащие углы равны;
\item внутренние односторонние в сумме дают $180^\circ$.
\end{itemize}

\section*{\color{darkaccent}8. Окружность}
\[
d=2r.
\]
Хорда через центр является диаметром. Радиус в точку касания перпендикулярен касательной.
Центральный угол равен соответствующей дуге, вписанный равен половине дуги. Если они опираются на одну дугу:
\[
\boxed{\text{центральный угол}=2\cdot\text{вписанный угол}}.
\]
Через три точки, не лежащие на одной прямой, проходит ровно одна окружность.

\section*{\color{darkaccent}9. Статистика}
Для значений $x_i$ с частотами $n_i$:
\[
\bar x=\frac{x_1n_1+\cdots+x_kn_k}{n_1+\cdots+n_k}.
\]
Медиана находится только после упорядочивания данных. При чётном числе значений берётся среднее двух центральных. Мода --- самое частое значение. Размах:
\[
R=x_{\max}-x_{\min}.
\]
Для круговой диаграммы:
\[
\alpha=360^\circ\cdot\frac{\text{часть}}{\text{целое}}.
\]

\section*{\color{darkaccent}10. Типичные ошибки}
\begin{itemize}
\item Проверить только верхнюю границу в неравенстве треугольника.
\item Принять любую медиану за высоту или биссектрису.
\item Применить свойство медианы к гипотенузе вне прямоугольного треугольника.
\item Не упорядочить данные перед поиском медианы.
\item Посчитать среднее без учёта частот.
\item Считать равенство углов «по рисунку» без теоремы.
\item После нахождения $x$ забыть вернуться к тому, что требуется в вопросе.
\end{itemize}

\clearpage
\section*{\color{darkaccent}Тренировочный блок · Путь А}
\begin{enumerate}
\item Смежные углы относятся как $2:7$. Найдите оба угла.
\item Две стороны треугольника равны $8$ и $13$. Сколько положительных целых значений может принимать третья сторона?
\item Медиана к гипотенузе прямоугольного треугольника равна $8$ см. Найдите гипотенузу.
\item Оценки $2,3,4,5$ получили соответственно $2,3,5,2$ ученика. Найдите среднюю оценку, медиану, моду и размах.
\item В равнобедренном треугольнике $ABC$: $AB=BC$, $\angle ABC=120^\circ$. $BM$ --- медиана. На луче $BM$ за $M$ выбрана $F$, $AF\perp AB$, $MF=54$. Найдите $AB$.
\end{enumerate}

\clearpage
\section*{\color{darkaccent}Ответы}
\begin{tabularx}{\linewidth}{@{}cX@{}}
\toprule
№ & Ответ\\
\midrule
1 & $40^\circ$ и $140^\circ$\\
2 & $15$ значений: $6,7,\ldots,20$\\
3 & $16$ см\\
4 & $\bar x=\frac{43}{12}=3\frac{7}{12}$; медиана $4$; мода $4$; размах $3$\\
5 & $36$\\
\bottomrule
\end{tabularx}
\end{document}